\subsection{The Alternating Group}

\begin{definition}
    Two cycles are called transpositions
\end{definition}

\begin{definition}
    An element $\sigma\in S_n$ is called even(odd)
    if it can be written as the product of an even(odd) number
    of transpositions
\end{definition}

\begin{definition}
    The alternating group, $A_n$, is the subgroup
    of $S_n$ made up of all even permutations.
    \[A_n=\left\{ \sigma \in S_n \mid \sigma \t{ is even} \right\}\]
\end{definition}

\begin{proposition}
    Every $\sigma \in S_n$ can be written as the product of transpositions.
\end{proposition}

\begin{proof}
    Since $\sigma$ can be written as the product disjoint cycles we only have to
    show that one cycles can be written as the product of transpositions.

    Observe $(a_1,a_2,\dots,a_m)$=$(a_1,a_m)(a_1,a_{m-1})\dots(a_1,a_2)$

    A $m$ cycle breaks into $m-1$ transpositions, so an even length cycle is odd 
\end{proof}

\begin{proposition}
    If $(a_1,b_1)(a_2,b_2)\dots(a_n,b_n)=(1)$ then $n$ is even, and thus $(1)\in A_n$
\end{proposition}

Observations:
\begin{align}
    (bc)(ab)&=(ac)(bc)\\
    (ac)(ab)&=(ab)(bc)\\
\end{align}
\begin{proof}
    Base Case: $n=1 \implies (ab)\neq (1)$
    
    $n=2 \implies (ab)(cd)=(1) \implies (ab)=(cd)$

    I.H. Assume for all $1 \leq k < m$ we have if $(a_1b_1)\dots(a_kb_k)=(1)$

    Assume $(a_1,b_1)(a_2,b_2)\dots(a_n,b_n)=(1)$

    Set $a_m=a$ and use $(1)$ or $(2)$ to move the rightmost appearance of $a$ to the left until

    1. the two cycle containing $a$ is next to its inverse
    
    $(a_1b_1)\dots(ab_m)=(a_1b_1)\dots(a_{j-1}b_{j-1})(ab)(ab)(a_{j+2}^`b_{j+1}^`)\dots(a_m^`b_m^`)$
    
    $=$ a product of $m-2$ transpositions by I.H. we have $m-2$ is even $\implies m$ is even

    2. the two cycle containing $a$ never finds an inverse

    $(a_1b_1)\dots(ab_m)=(ab)(a_2^`b_2^`)\dots(a_m^`b_m^`) \implies a$ is not fixed $\implies \neq (1)$
\end{proof}
\newpage
\begin{proposition}
    $A_n$ is a subgroup of order $\frac{n!}{2}$
\end{proposition}

\begin{proof}
    $(1)\in A_n$ so nonempty
    Assume $\sigma,\mu \in A_n \implies \sigma=a_1a_2\dots a_r$ and $\mu=b_1b_2\dots b_s$
    where $r$ and $s$ are even

    $\sigma\mu^{-1}=a_1\dots a_rb_s\dotsb_1$ $r+s$ is even so $\sigma\mu^{-1}\in A_n$

    Thus, $A_n$ is a subgroup

    Even = odd $\implies |A_n|=\frac{n!}{2}$
\end{proof}

\begin{corollary}
    $[S_n:A_n]=2$ since \#even = \#odd 
    $\implies A_n \unlhd S_n \implies S_n\diagup A_n \cong \Z_2$
\end{corollary}

\begin{lemma}
    \label[lemma]{lem:alternating3cycle}
    For $n \geq 3$, $A_n$ is generated by 3-cycles.
\end{lemma}

\begin{proof}
    Take $\sigma \in S_n$, write $\sigma=a_1b_2\dots a_mb_m$

    $a_jb_j=\begin{cases} 
        (ab)(cd)=(acb)(acd)\\
        (ab)(ac)=(acb) 
     \end{cases}$
    $\implies \sigma$ can be written as 3-cycles
\end{proof}

\begin{lemma}
    \label[lemma]{lem:noramlsubgrouptoalternating3cycle}
    Let $N \unlhd A_n$ with $n \geq 3$. If $N$ contains a 3-cycle
    then $N = A_n$
\end{lemma}
\begin{proof}
    Note all 3-cycles can be written as products of $(12\_)$. 
    $a,b,c\in \{3,\dots n\}$
    \begin{align*}
        (1a2)&=(12a)(12a)\\
        (1ab)&=(12b)(12a)^2\\
        (2ab)&=(12b)^2(12a)\\
        (abc)&=(12a)^2(12c)(12b)^2(12a)
    \end{align*}

    Take $(abc)\in N$

    $(2b)(1a)(abc)(1a)(2b)=(12c) \in N$

    Take $m \in \{3,\dots,n\}, (12m)=(12)(mc)(12c)^2(mc)(12)\in N \implies$
    all 3-cycles are in $N$, but 3-cycles generate $A_n$ so $N=A_n$
\end{proof}
\newpage
\begin{theorem}
    \label[theorem]{thm:alternatingnormalsubgroup}
    For $n \geq 5$, if $N \unlhd A_n$ then $N=\{e\}$ or $N=A_n$
\end{theorem}

\begin{proof}
    Take $(1)\neq \sigma \in N$ and write it as a product of 
    disjoint cycles.

    Case 1: There is a cycle of length $\geq 4$

    Write: $\sigma=(a_1a_2a_3a_4\dots a_m)\mu$

    $(a_1a_3a_m)=\sigma^{-1}((a_1a_2a_3)\sigma(a_1a_2a_3)^{-1})$ 
    since $N$ is normal 
    $\implies \in N \implies N = A_n$ by \Cref{lem:alternating3cycle}

    Case 2: $\sigma$ only contains 2,3-cycles

    Case 2a: $\sigma$ has at least two 3-cycles
    
    $\sigma=(a_1a_2a_3)(a_4a_5a_6)\mu 
    \implies (a_1a_4a_2a_6a_3)=\sigma^{-1}((a_1a_2a_4)\sigma(a_1a_2a_4)^{-1})\in N$

    $\implies$ Case 1 $\implies N = A_n$

    Case 2b: $\sigma$ has one 3-cycle

    $\sigma=(a_1a_2a_3)b_2b_3\dots b_k$

    $\sigma^2=(132) \in N \implies N=A_n$

    Case 3: $\sigma$ has only 2-cycles

    $\sigma=(a_1a_2)(a_3a_4)\mu$

    $(a_1a_3)(a_2a_4)=\sigma^{-1}((a_1a_2a_3)\sigma(a_1a_2a_3)^{-1}) \in N$

    $(a_1a_3a_5)=(a_1a_3)(a_2a_4)((a_1a_3a_5)(a_1a_3)(a_2a_4)(a_1a_3a_5)^{-1}) \in N
    \implies N = A_n$

\end{proof}

\begin{definition}
    We say a group $G$ is simple if it has no proper normal subgroup
\end{definition}

\newpage
\subsubsection{Exercises}
\begin{example}
    $\Z_n$ is simple $\iff n$ is prime
    \begin{proof}
        Assume $\Z_n$ is simple

        Since all subgroups of $\Z_n$ are normal $\implies$ the only subgroups
        of $\Z_n$ are $\Z_n \t{ and } \{e\}$

        By \Cref{thm:lagrange} only $1,n \mid n \implies n$ is prime

        Assume $n$ is prime

        Since all subgroups of $\Z_n$ are normal $\implies$ the only subgroups
        of $\Z_n$ are $\Z_n \t{ and } \{e\}$ since subgroups order have to divide $n$

        Thus, $\Z_n$ is simple
    \end{proof}
\end{example}